% Version 3.5: insert in the Cross-region Handler subsection, before Protocol 3.
\paragraph{Complete source inputs and ordered replay.}
\label{sec:source-step}
Each correct regional member retains its authenticated source cursors,
verified interpreter checkpoint and public dependency index across
session replacement. Source availability invokes the handler below;
it is not an instruction to enqueue an independently scheduled replay
job. The handler uses the existing complete-message processing bound
and atomic local state-machine steps. It does not change the canonical
order of \cref{alg:force,alg:external}.

\paragraph{A common record-completeness rule.}
\label{sec:record-completeness}
For an EXPORT at position $p$, its public dependency basis consists of
the authenticated source prefix strictly before $p$ and the protocol's
authorized checkpoint identified by that canonical source context.
A verifier's choice of a newer local starting checkpoint does not change
this basis. The designated checkpoint specifies
the commitments and source identities of its dependency objects. An
EXPORT must supply the bytes of each transitive validation dependency,
or identify those same bytes by their exact commitment and position in
this public basis. The referenced object must actually occur there;
a bare hash, URL, future authority name or sender assertion is not a
public dependency. Source authentication does not itself establish
regional or application validity: those checks still apply to every
supplied or referenced object. The DATA reference bounds remain
unchanged, including strict precedence before its containing Espresso
block.

This rule depends on the record and its canonical public basis, not
on a verifier's private cache. A cache can avoid repeated retrieval or
validation of the same authenticated object. A private cache hit cannot
make an otherwise incomplete EXPORT valid. Conversely, a local cache
miss cannot make a publicly complete EXPORT invalid: retrieve its
specified public object and resume at the same position. Two verifiers
with different caches therefore apply the same completeness predicate.

An incomplete namespace download, incomplete applicable control prefix,
or unavailable authority genuinely required to resolve the current
canonical boundary pauses the source cursor. None is treated as empty
history or an ineligible region. By contrast, within a complete
namespace, a candidate that omits required nonpublic proof bytes or
relies only on a reference outside its public basis fails record completeness.
Reject that record without changing $G$ or its counts and continue in
namespace order; do not fetch arbitrary future material on its behalf.
Check the domain, current base, open version and structural append
permission before undertaking its deeper dependency checks. A closed-beta
record may retain useful historical evidence, but its optional archival
validation cannot block current public replay.

\paragraph{The regional source packet.}
A \tagg{SOURCESTEP} is an authenticated message to a fixed regional
snapshot. It contains the following fields.
\begin{table}[htbp]\centering\small
\begin{tabularx}{\linewidth}{@{}>{\raggedright\arraybackslash}p{.25\linewidth}Y@{}}\toprule
Field & Required value\\\midrule
Domain and routing & Chain and protocol domain, encoding version,
sender identity and the recipient snapshot's immutable authority
identity. Routing does not itself select a session.\\
Source position & Espresso height $q$, full header $J_q$, finality
and linkage evidence, complete relevant namespace bytes or authenticated
absence, and the canonical record order.\\
Control context & The complete applicable finalized L1 control context,
including pending records and its completeness evidence; the authority
objects required by the current boundary, when available.\\
Predecessor basis & The preceding source position and the identifier of
an authorized or already verified replay checkpoint. The recipient
derives subsequent interpreter states itself; a sender-supplied new
checkpoint is not trusted.\\
Dependencies & Full bytes and source evidence for required public
references not already resolved by that checkpoint, with a normalized
index linking each object to its commitment and canonical source
position. Candidate-local proof bytes remain part of the namespace
record and obey \cref{sec:record-completeness}.\\\bottomrule
\end{tabularx}
\end{table}

On acquiring and authenticating a new complete source input, every
correct member sends its source packet to its installed regional
snapshot in that same availability event, before waiting for its local
P6 replay result. It need not know which region that replay will select.
It includes the available authenticated boundary dependencies; a genuine
missing boundary authority remains an explicit wait, not a fabricated
complete step. Native preparation is separate and never delays this
source relay. Equal source packets are deduplicated by source position
and commitment, without suppressing the first correct relay.

\paragraph{Receive, drain and activate.}
The receiver authenticates the packet and its complete dependencies,
retains it by source position, and processes the contiguous ready prefix
from its maintained checkpoint. An out-of-order packet is buffered;
arrival of its missing predecessor immediately resumes this same
ordered handler. There is no additional polling tick or replay queue.
For each ready block, the handler executes P6, including the P5 prelude,
the common record-completeness gate and EXPORT validation in order.
It durably records the resulting cursor, eligibility, duty, window,
boundary and branch state before exposing them. A local receive time
never replaces a canonical opening timestamp.

The same event delivers the verified transition to the Session Manager.
The Manager fixes the derived handoff, checks the current activation
predicate and, when the exact executable anchor state and remaining
Ready dependencies are already prepared, performs Ready and initial
Enter atomically under \cref{alg:bootstrap}. It tests and sets the
durable entry flag once for this pair. It neither reruns Force nor
increments beta itself. Duplicate inputs cannot reset a window, reuse
old locks or votes, reset entry times, or enter a pair twice. The new
forced version still cannot export in its triggering block.

On closure, atomically persist the authenticated version-wide
\tagg{Closed} gate and its durable control-routing obligations.
Every initial entry, fresh core action and delayed callback consults
this gate by its own beta, regardless of program phase. No scan of old
workers, archive copy or completion of retirement is a prerequisite for
gating the old version or entering a valid new pair. Version-indexed
control records retain the obligations for all known snapshots;
missing snapshot context remains pending. This organization does not
defer a complete-close handler or multicast required by
\cref{sec:retirement}. Whenever a complete old-session context is
received, including at a Waiting, Completed or Terminated record, that
receive handler validates it, completes and relays closure before its
ordinary response, under the existing regional complete-message bound.
The original per-session $c_v$, $4\Delta$ drain rule and evidence duties
remain unchanged. Recording a Closed gate alone does not start that
drain clock or authorize an early exit.

\paragraph{Prepared-tail timing.}
\label{sec:source-step-timing}
Consider a tail with fixed $G^*$ and installed slot $r^*$, after the
finite required authority prefix, anchor provenance and native
preparation have completed and are retained. Use the continuously
online, history-following regional participants covered by
\cref{ass:espresso-visibility,ass:network}; their prior verified replay
checkpoint is retained. Let $J_q$ open a retry for region $i$ and first
become final at real time $a$. Every predecessor $J_{q'}$, $q'\le q$,
became final no later than $a$. Its complete source evidence is therefore
available by $a+T_E$. Applicable controls are available before their
crossings, and the prepared boundary dependencies are cached. A
malformed candidate cannot create an additional external wait under
\cref{sec:record-completeness}.

The immediate regional sends occur by that common source deadline.
Their delivery and complete validation processing finish by
$a+T_E+\Delta_i$. The ordered handler advances the already maintained
checkpoint through the available inputs using atomic logical steps;
it does not insert another network round at every predecessor height.
Its Manager step then attaches the already executable anchor object
and performs initial Enter. Thus, if the target remains activatable,
the first correct initial entry satisfies
\[
 e(\beta,r,0,0)\le a+T_E+\Delta_i.
\]
If a later verified closure has already intervened, the closure
alternative applies instead. Ready needs the opening header, not a
header generated after it. This argument uses no block-generation rate
or lower-lag bound; growing Espresso time and the upper-lead condition
retain their separate roles in the progress proof.

The timing statement concerns complete message validation and logical
state attachment. A persistent executable anchor object can be attached
to a fresh continuation without importing old consensus state; a
serialized snapshot that still requires database import, copying or
native-runtime initialization does not satisfy this prepared contract.
Those external operations must finish before the bound is invoked.
Nor does the argument give a constant bound for a returning replica's
arbitrary old backlog. A concrete native adapter must establish that it
implements this interface; its performance is not proved merely by
naming the operation a handler. This packet adds neither a new service
assumption nor a READY certificate or Espresso publication obligation.
